%======================================================================
\section{Attributing the spread}\label{sec:attribution}
%======================================================================

The market quotes AT1 bonds by their yield to the first call over the matching
Treasury. We want to say how much of that spread pays for each risk. This section
defines the attribution and discusses its dependence on the order in which risks are
switched on.

\subsection{Features, prices and yields}

Index the risk features of the bond by the set
\[
  \mathcal N=\{\mathrm R,\mathrm S,\mathrm M,\mathrm C,\mathrm E\},
\]
standing for resolution hazard ($\lres$), AT1-specific hazard ($\lspec$), MDA coupon
restriction, conversion at the trigger, and extension (the economic call rule
\eqref{eq:call} instead of a certain call at $T_1$). For any subset $K\subseteq\mathcal N$
let $V(K)$ be the dirty price with exactly the features in $K$ switched on, computed on
the fixed set of paths:
\begin{itemize}
  \item hazard $(\lres\1_{\mathrm R\in K}+\lspec\1_{\mathrm S\in K})\,s^+_t/s_0$;
  \item coupons scaled by $\phi$ if $\mathrm M\in K$, paid in full otherwise;
  \item conversion at $\tau_m$ if $\mathrm C\in K$, no trigger otherwise;
  \item call by \eqref{eq:call} if $\mathrm E\in K$, certain call at $T_1$ otherwise.
\end{itemize}
$V(\emptyset)$ is the price of a risk-free bullet to $T_1$, and $V(\mathcal N)$ is the
model price, equal to the market price by calibration.

Let $y(V)$ be the semi-annual yield to the first call implied by a dirty price $V$,
treating the bond as a bullet to $T_1$, and define the spread of a subset as
$v(K)=y(V(K))-y(V(\emptyset))$. Because $V(\emptyset)$ is discounted at the risk-free
rate, $y(V(\emptyset))$ is the Treasury yield and $v(\mathcal N)$ is the market
yield-to-call spread. The map $V\mapsto y(V)$ is strictly decreasing, so $v$ is
increasing in risk.

\subsection{Sequential attribution and its order dependence}

Given an order $\pi=(\pi_1,\dots,\pi_5)$ of the features, the sequential attribution is
\begin{equation}\label{eq:sequential}
  \varphi^\pi_{\pi_k}=v(\{\pi_1,\dots,\pi_k\})-v(\{\pi_1,\dots,\pi_{k-1}\}),
\end{equation}
which telescopes to $\sum_k\varphi^\pi_{\pi_k}=v(\mathcal N)$: the components add up
exactly to the market spread. The attribution depends on $\pi$ whenever the features
interact, that is whenever $v$ is not additive. Interactions are pervasive here. The MDA
cuts coupons that would otherwise be paid only if the bond survives, so its cost depends on
the hazard; the value of the call option depends on the discount rate of the extended
cash flows, so its cost depends on the hazard and on the MDA; and the yield transformation
is non-linear.

The base attribution of the short note uses the order $(\mathrm R,\mathrm S,\mathrm M,\mathrm C,\mathrm E)$:
the hazard first, so that every contract feature is valued with risky discounting, and
extension last, so that the call is valued with every other feature in place.

\subsection{Shapley values}

The Shapley value \citep{Shapley1953} averages the sequential attribution over all
$|\mathcal N|!=120$ orders:
\begin{equation}\label{eq:shapley}
  \varphi^{\mathrm{Sh}}_j=\sum_{K\subseteq\mathcal N\setminus\{j\}}
  \frac{|K|!\,(|\mathcal N|-|K|-1)!}{|\mathcal N|!}\bigl(v(K\cup\{j\})-v(K)\bigr).
\end{equation}
It is the unique attribution satisfying efficiency ($\sum_j\varphi_j=v(\mathcal N)$),
symmetry (interchangeable features receive the same value), the null-player property
and additivity. It needs the $2^5=32$ values $v(K)$, which on a fixed set of paths require
only eight cash-flow profiles per bond (one for each combination of M, C and E): the two
hazard features enter through the discount exponent and are evaluated analytically.

The full Shapley value treats all orders as equally meaningful. For AT1 bonds they are
not. Orders in which extension is switched on before any hazard value the issuer's option
on a risk-free bond. When a risk-free bond pays its reset margin $m$ over Treasuries, not
calling it is a gift to the investor, so these orders assign extension a large negative
value; this is precisely the inconsistency that the state-dependent hazard of
Section~\ref{sec:hazard} is designed to avoid. We therefore also compute a Shapley value
within blocks.

\begin{definition}[Nested Shapley value]
Partition the features into the hazard block $\mathcal H=\{\mathrm R,\mathrm S\}$ and the
contract block $\mathcal K=\{\mathrm M,\mathrm C,\mathrm E\}$. The nested value is the
Shapley value of the game $v$ restricted to $\mathcal H$ for $j\in\mathcal H$, and the
Shapley value of the game $K\mapsto v(\mathcal H\cup K)-v(\mathcal H)$ on $\mathcal K$ for
$j\in\mathcal K$.
\end{definition}

The nested value is symmetric within each block and enters the hazard block first. It is
related to the value for games with a priori unions of \citet{Owen1977}, which would also
average over the order of the two blocks; we fix the order of the blocks for the economic
reason just given. By construction it is efficient: the hazard block receives
$v(\mathcal H)$ and the contract block $v(\mathcal N)-v(\mathcal H)$.

\begin{proposition}[Split of the hazard block]\label{prop:hazardsplit}
Within the hazard block, the nested value is
\[
  \varphi_{\mathrm R}=\tfrac12\bigl[v(\{\mathrm R\})+v(\mathcal H)-v(\{\mathrm S\})\bigr],\qquad
  \varphi_{\mathrm S}=\tfrac12\bigl[v(\{\mathrm S\})+v(\mathcal H)-v(\{\mathrm R\})\bigr].
\]
If the spread were linear in the hazard level on the relevant range,
$v(\{\mathrm R\})+v(\{\mathrm S\})=v(\mathcal H)$ and the nested value would coincide with
the sequential attribution in either order, giving
$\varphi_{\mathrm R}/(\varphi_{\mathrm R}+\varphi_{\mathrm S})=\lres/\lambda$.
\end{proposition}

\begin{proof}
The first statement is \eqref{eq:shapley} for two players. Under linearity
$v(\{\mathrm R\})=\lres\,g$ and $v(\{\mathrm S\})=\lspec\,g$ for some slope $g$, and both
orders give the same increments.
\end{proof}

For BNP the yield-to-call spread is very close to linear in the hazard on the relevant
range, so the split of the hazard block is effectively order-free: the nested value and
the base order agree to within $0.2$bp (Section~\ref{sec:robust}). The order matters for the
contract block, and there the attribution should be reported as a range.
